% =========================================================
% TABLA 3: OBJETIVO ESPECÍFICO 3
% =========================================================
\begin{table}[htbp]
	\centering
	\caption{Cumplimiento del Objetivo Específico 3}
	\label{tab:cumplimiento-objetivo-3}
	\begin{tabular}{|>{\raggedright\arraybackslash}p{0.43\textwidth}|>{\raggedright\arraybackslash}p{0.51\textwidth}|}
		\hline
		\rowcolor[gray]{0.85}
		\multicolumn{2}{|c|}{\textbf{\textit{Objetivo específico 3}}}                                                                                                                                                                                                                                                                                                                                                                                     \\
		\hline
		\multicolumn{2}{|>{\centering\arraybackslash}p{\dimexpr0.94\textwidth+2\tabcolsep+\arrayrulewidth\relax}|}{\textit{Seleccionar una alternativa de software libre/código abierto para la implementación de un servicio de directorio en la infraestructura del grupo de investigación GRID}}                                                                                                                                                                    \\
		\hline
		\textbf{\textit{Detalle de cumplimiento}}                                                                                                                                                                                                                                                                                                                                                                                    & \textbf{\textit{Observaciones}} \\
		\hline
		Se aplicó la metodología \textit{\textcolor[HTML]{1B6B93}{\ac{DAR}}} del marco \ac{CMMI} sobre las cinco tecnologías candidatas, a partir de criterios de evaluación derivados de los requisitos definidos (ver capítulo \textcolor[HTML]{1B6B93}{\textbf{\ref{cap:dar}}}). Como resultado de la evaluación estructurada, se seleccionó \textit{\textcolor[HTML]{1B6B93}{FreeIPA}} como la tecnología base para la solución. &
		\textbf{Paso 1: Definición de los criterios de evaluación.}\par\vspace{0.4em}
		- Se definieron trece criterios derivados de los requisitos y de aspectos de gestión de identidades y seguridad, ponderados según su importancia para el proyecto (sección \textcolor[HTML]{1B6B93}{\ref{sec:dar-metodologia}} y sección \textcolor[HTML]{1B6B93}{\ref{subsec:dar-lista-criterios}}).\par\vspace{0.6em}
		\textbf{Paso 2: Ponderación y valoración de las alternativas.}\par\vspace{0.4em}
		- Se asignó la ponderación a cada criterio y se valoraron las cinco alternativas en una escala de 0 a 5 (sección \textcolor[HTML]{1B6B93}{\ref{subsec:dar-ponderacion}}).\par\vspace{0.6em}
		\textbf{Paso 3: Análisis \textit{\textcolor[HTML]{1B6B93}{DAR}} y selección.}\par\vspace{0.4em}
		- \textit{FreeIPA} obtuvo la mayor puntuación ponderada global, destacándose por su gestión centralizada de identidades, políticas de contraseñas e infraestructura de certificados integrada (sección \textcolor[HTML]{1B6B93}{\ref{subsec:dar-analisis}}, resultados en \autoref{fig:dar} y sección \textcolor[HTML]{1B6B93}{\ref{sec:dar-resultados}}).                                                                                                     \\
		\hline
	\end{tabular}
\end{table}
